% ===== PRÉAMBULE CANONIQUE — à coller VERBATIM en tête de CHAQUE .tex =====
\documentclass[11pt,a4paper]{article}
\usepackage[utf8]{inputenc}\usepackage[T1]{fontenc}\usepackage{lmodern}
\usepackage[french]{babel}
\usepackage{amsmath,amssymb,mathtools}\usepackage{icomma}
\usepackage[version=4]{mhchem}          % \ce{...} : équations & formules
\usepackage{siunitx}\sisetup{locale=FR} % \qty{}{}, \si{}, \ang{}
\usepackage{chemfig}                    % \chemfig{...} : structures développées
\usepackage{xcolor}\usepackage{colortbl}
\usepackage{tikz}\usepackage{pgfplots}\pgfplotsset{compat=1.18}
\usepackage[most]{tcolorbox}\usepackage{enumitem}\usepackage{array,booktabs}
\usepackage[a4paper,margin=2.2cm]{geometry}
\usetikzlibrary{arrows.meta,calc,angles,quotes,positioning,patterns,backgrounds,shapes.geometric}
\definecolor{primary}{RGB}{222,49,99}\definecolor{primarydark}{RGB}{150,22,55}
\definecolor{defcol}{RGB}{0,114,178}\definecolor{methcol}{RGB}{52,92,148}
\definecolor{excol}{RGB}{0,135,90}\definecolor{attncol}{RGB}{200,40,55}
\tcbset{boxbase/.style={enhanced,breakable,boxrule=0.9pt,arc=2.5pt,
  left=3mm,right=3mm,top=2.3mm,bottom=2.3mm,fonttitle=\bfseries,coltitle=white}}
% --- boîtes TUTORIEL (noms EXACTS, ne pas renommer) ---
\newtcolorbox{cle}[1][]{boxbase,colback=defcol!6,colframe=defcol,
  title={\textsc{À retenir}\ifx\relax#1\relax\relax\else\textmd{ : \itshape #1}\fi}}
\newtcolorbox{methode}[1][]{boxbase,colback=methcol!7,colframe=methcol,sharp corners,
  title={\textsc{Méthode}\ifx\relax#1\relax\relax\else\textmd{ --- \itshape #1}\fi}}
\newtcolorbox{exemplebox}[1][]{boxbase,colback=excol!5,colframe=excol,
  title={\textsc{Exemple}\ifx\relax#1\relax\relax\else\textmd{ : \itshape #1}\fi}}
\newtcolorbox{piege}[1][]{boxbase,colback=attncol!6,colframe=attncol,
  title={\textsc{Piège}\ifx\relax#1\relax\relax\else\textmd{ : \itshape #1}\fi}}
\newtcolorbox{remarque}[1][]{boxbase,colback=black!4,colframe=black!45,
  title={\textsc{Remarque}\ifx\relax#1\relax\relax\else\textmd{ : \itshape #1}\fi}}
% --- boîtes EXERCICE / CORRIGÉ (noms EXACTS, auto-comptées) ---
\newtcolorbox[auto counter]{exo}[1][]{boxbase,colback=primary!4,colframe=primary,
  title={\textsc{Exercice}~\thetcbcounter\ifx\relax#1\relax\relax\else\textmd{ --- \itshape #1}\fi}}
\newtcolorbox[auto counter]{corrige}[1][]{boxbase,colback=excol!4,colframe=excol,
  title={\textsc{Corrigé}~\thetcbcounter\ifx\relax#1\relax\relax\else\textmd{ --- \itshape #1}\fi}}
\newcommand{\R}{\mathbb{R}}\newcommand{\N}{\mathbb{N}}\newcommand{\Z}{\mathbb{Z}}\newcommand{\ds}{\displaystyle}
% (un \usepackage{fancyhdr}+en-tête est possible ; il est ignoré côté web)
% =========================================================================
\begin{document}
\begin{center}\color{primarydark}\large\bfseries Thème 3 — Potentiels standard et prévision \quad$\bullet$\quad Niveau Difficile\end{center}

\begin{exo}[vrai ou faux : qui oxyde qui ?]
Données : $E_{\mathrm{r}}^{\circ}$ (V) : \ce{Cl2}/\ce{Cl^-} $=+1{,}36$ ;
\ce{Fe^{3+}}/\ce{Fe^{2+}} $=+0{,}77$ ; \ce{I2}/\ce{I^-} $=+0{,}54$ ; \ce{Cu^{2+}}/\ce{Cu} $=+0{,}34$ ;
\ce{Zn^{2+}}/\ce{Zn} $=-0{,}76$. Vrai ou faux ? Combien de propositions sont vraies ?
\begin{enumerate}[label=\alph*)]
  \item Le dichlore \ce{Cl2} oxyde les ions \ce{Fe^{2+}} en \ce{Fe^{3+}}.
  \item Le diiode \ce{I2} oxyde les ions \ce{Fe^{2+}} en \ce{Fe^{3+}}.
  \item Le cuivre métallique réduit les ions \ce{Zn^{2+}} en \ce{Zn}.
  \item Les ions \ce{Fe^{3+}} oxydent le cuivre métallique.
\end{enumerate}
\end{exo}

\begin{exo}[ces réactions ont-elles lieu ?]
Dans les conditions standard, dis si chaque réaction est possible. Données :
$E_{\mathrm{r}}^{\circ}(\ce{Cl2}/\ce{Cl^-})=+1{,}36$ ; $E_{\mathrm{r}}^{\circ}(\ce{Br2}/\ce{Br^-})
=+1{,}06$ ; $E_{\mathrm{r}}^{\circ}(\ce{I2}/\ce{I^-})=+0{,}54$ ;
$E_{\mathrm{r}}^{\circ}(\ce{Ag+}/\ce{Ag})=+0{,}80$ ; $E_{\mathrm{r}}^{\circ}(\ce{Cu^{2+}}/\ce{Cu})
=+0{,}34$~V.
\begin{enumerate}[label=\alph*)]
  \item \ce{2 I^- + Br2 -> I2 + 2 Br^-}
  \item \ce{2 Cl^- + I2 -> Cl2 + 2 I^-}
  \item \ce{Cu + 2 Ag+ -> Cu^{2+} + 2 Ag}
\end{enumerate}
\end{exo}

\begin{exo}[le difluor dissout-il l'or ?]
On plonge une pièce d'or dans une solution saturée de difluor \ce{F2}. Données :
$E_{\mathrm{r}}^{\circ}(\ce{Au^{3+}}/\ce{Au})=+1{,}52$~V ;
$E_{\mathrm{r}}^{\circ}(\ce{F2}/\ce{F^-})=+2{,}87$~V.
\begin{enumerate}[label=\alph*)]
  \item L'or se dissout-il ? Justifie.
  \item Sachant que \ce{Cl2} ($+1{,}36$~V) et \ce{O2} ($+1{,}23$~V) sont les oxydants courants les plus
        forts après \ce{F2}, expliquent pourquoi ils ne dissolvent pas l'or.
\end{enumerate}
\end{exo}

\begin{exo}[médiamutation du manganèse]
Le manganèse intervient dans trois couples : $E_{\mathrm{r}}^{\circ}(\ce{MnO4^-}/\ce{MnO2})
=+1{,}69$~V ; $E_{\mathrm{r}}^{\circ}(\ce{MnO2}/\ce{Mn^{2+}})=+1{,}22$~V ;
$E_{\mathrm{r}}^{\circ}(\ce{MnO4^-}/\ce{Mn^{2+}})=+1{,}51$~V. En milieu acide, l'ion permanganate
\ce{MnO4^-} et l'ion \ce{Mn^{2+}} peuvent-ils réagir ensemble pour former \ce{MnO2}
(médiamutation) ? Identifie les deux couples mis en jeu et conclus par $\Delta E^{\circ}$.
\end{exo}

\end{document}
